% paper27-certificate-chains.tex
% Certificate Chains: Compositional Proof-Carrying Code Architecture
% Paper 27 of the Judgment Geometry series.
% Compile: pdflatex paper27-certificate-chains
\documentclass[11pt]{article}
% jugeo-common.tex — shared preamble for all Judgment Geometry papers
% Include via: % jugeo-common.tex — shared preamble for all Judgment Geometry papers
% Include via: % jugeo-common.tex — shared preamble for all Judgment Geometry papers
% Include via: \input{jugeo-common}

% ── Packages ────────────────────────────────────────────────────────────
\usepackage{amsmath,amssymb,amsthm,mathtools}
\usepackage{tikz-cd}
\usepackage{listings}
\usepackage{xcolor}
\usepackage{xspace}
\usepackage{booktabs}
\usepackage{multirow}
\usepackage{enumitem}
\usepackage[expansion=false]{microtype}
\usepackage{hyperref}
\usepackage{cleveref}
% \usepackage{thmtools}  % disabled: not available in this TeX installation
\usepackage{float}
\usepackage{caption}
\usepackage{subcaption}
\usepackage{tabularx}
\usepackage{stmaryrd}       % \llbracket, \rrbracket
\usepackage{bbm}            % \mathbbm{1}

% ── Page geometry ───────────────────────────────────────────────────────
\usepackage[margin=1in]{geometry}

% ── Theorem environments ────────────────────────────────────────────────
\theoremstyle{plain}
\newtheorem{theorem}{Theorem}[section]
\newtheorem{lemma}[theorem]{Lemma}
\newtheorem{proposition}[theorem]{Proposition}
\newtheorem{corollary}[theorem]{Corollary}

\theoremstyle{definition}
\newtheorem{definition}[theorem]{Definition}
\newtheorem{example}[theorem]{Example}
\newtheorem{construction}[theorem]{Construction}

\theoremstyle{remark}
\newtheorem{remark}[theorem]{Remark}
\newtheorem{notation}[theorem]{Notation}

% ── Python listings style ──────────────────────────────────────────────
\definecolor{jugeoBlue}{HTML}{2563EB}
\definecolor{jugeoGreen}{HTML}{059669}
\definecolor{jugeoRed}{HTML}{DC2626}
\definecolor{jugeoPurple}{HTML}{7C3AED}
\definecolor{jugeoGray}{HTML}{6B7280}
\definecolor{jugeoBg}{HTML}{F8FAFC}

\lstdefinestyle{jugeo-python}{
  language=Python,
  basicstyle=\ttfamily\small,
  keywordstyle=\color{jugeoBlue}\bfseries,
  stringstyle=\color{jugeoGreen},
  commentstyle=\color{jugeoGray}\itshape,
  numberstyle=\tiny\color{jugeoGray},
  backgroundcolor=\color{jugeoBg},
  numbers=left,
  numbersep=6pt,
  frame=single,
  framerule=0.4pt,
  rulecolor=\color{jugeoGray!40},
  breaklines=true,
  breakatwhitespace=true,
  showstringspaces=false,
  tabsize=4,
  xleftmargin=1.5em,
  framexleftmargin=1.5em,
  morekeywords={dataclass,frozen,field,Enum,IntEnum,auto,Any,
                Mapping,Sequence,Optional,match,case},
  literate={→}{{$\to$}}1 {←}{{$\leftarrow$}}1
           {≤}{{$\leq$}}1 {≥}{{$\geq$}}1 {≠}{{$\neq$}}1
           {∈}{{$\in$}}1 {∀}{{$\forall$}}1 {∃}{{$\exists$}}1
           {⊕}{{$\oplus$}}1 {⊖}{{$\ominus$}}1,
  escapeinside={(*@}{@*)},
}
\lstset{style=jugeo-python}

\lstdefinestyle{jugeo-output}{
  basicstyle=\ttfamily\small,
  backgroundcolor=\color{jugeoBg},
  frame=single,
  framerule=0.4pt,
  rulecolor=\color{jugeoGray!40},
  breaklines=true,
  showstringspaces=false,
  xleftmargin=1.5em,
  framexleftmargin=0.5em,
  numbers=none,
}

% ── JuGeo-specific macros ──────────────────────────────────────────────

% Judgment 8-tuple
\newcommand{\J}{\mathcal{J}}
\newcommand{\Jud}[8]{(#1,\,#2,\,#3,\,#4,\,#5,\,#6,\,#7,\,#8)}
\newcommand{\JudShort}[1]{\J_{#1}}

% Components
\newcommand{\coord}{\mathsf{c}}            % coordinate
\newcommand{\prop}{\varphi}                 % proposition
\newcommand{\carrier}{\mathcal{A}}          % carrier type
\newcommand{\evid}{\mathcal{E}}             % evidence bundle
\newcommand{\oblig}{\mathcal{O}}            % obligations
\newcommand{\obstr}{\mathcal{B}}            % obstructions
\newcommand{\trust}{\mathcal{T}}            % trust annotation
\newcommand{\prov}{\Pi}                     % provenance

% Trust algebra
\newcommand{\Talg}{\mathfrak{T}}            % trust ordered algebra
\newcommand{\Eadm}{\mathcal{E}_{\mathrm{adm}}}
\newcommand{\tleq}{\preceq}                 % trust ordering
\newcommand{\tjoin}{\oplus}                 % conservative join
\newcommand{\tatten}{\ominus}               % attenuation
\newcommand{\tpromote}{\uparrow_\pi}        % promotion
\newcommand{\tdemote}{\downarrow_\chi}       % demotion

% Trust tiers
\newcommand{\Tcontra}{\ensuremath{\mathsf{CONTRADICTED}}}
\newcommand{\Tunver}{\ensuremath{\mathsf{UNVERIFIED}}}
\newcommand{\Tcopilot}{\ensuremath{\mathsf{COPILOT}}}
\newcommand{\Toracle}{\ensuremath{\mathsf{ORACLE}}}
\newcommand{\Truntime}{\ensuremath{\mathsf{RUNTIME}}}
\newcommand{\Tsolver}{\ensuremath{\mathsf{SOLVER}}}
\newcommand{\Tproof}{\ensuremath{\mathsf{PROOF}}}

% Site and geometry
\newcommand{\Site}{\mathcal{S}}             % semantic site
\newcommand{\Cov}{\mathrm{Cov}}            % covering family
\newcommand{\Ob}{\mathrm{Ob}}              % objects
\newcommand{\Mor}{\mathrm{Mor}}            % morphisms
\newcommand{\res}{\mathrm{res}}            % restriction
\newcommand{\Cech}{\check{C}}              % Čech complex
\newcommand{\Hcoh}[1]{H^{#1}}             % cohomology group

% Descent
\newcommand{\descent}{\mathrm{desc}}
\newcommand{\glue}{\mathrm{glue}}
\newcommand{\overlap}[2]{#1 \cap #2}

% Semantic moves
\newcommand{\VERIFY}{\textsc{Verify}}
\newcommand{\CONSTRUCT}{\textsc{Construct}}
\newcommand{\NORMALIZE}{\textsc{Normalize}}
\newcommand{\GLUE}{\textsc{Glue}}
\newcommand{\RESTRICT}{\textsc{Restrict}}
\newcommand{\TRANSPORT}{\textsc{Transport}}
\newcommand{\REPAIR}{\textsc{Repair}}
\newcommand{\PROMOTE}{\textsc{Promote}}
\newcommand{\CHALLENGE}{\textsc{Challenge}}
\newcommand{\ESCALATE}{\textsc{Escalate}}

% Fragments
\newcommand{\QFLIA}{\texttt{QF\_LIA}}
\newcommand{\QFLRA}{\texttt{QF\_LRA}}
\newcommand{\QFBV}{\texttt{QF\_BV}}
\newcommand{\QFUF}{\texttt{QF\_UF}}

% Misc
\newcommand{\Python}{\textsc{Python}}
\newcommand{\JuGeo}{\textsc{JuGeo}}
\newcommand{\LEAN}{\textsc{Lean}\,4}
\newcommand{\Fstar}{F$^{\star}$}
\newcommand{\Coq}{\textsc{Coq}}
\newcommand{\Dafny}{\textsc{Dafny}}

% Common shortcuts
\newcommand{\ie}{i.e.\@\xspace}
\newcommand{\eg}{e.g.\@\xspace}
\newcommand{\cf}{cf.\@\xspace}
\newcommand{\wrt}{w.r.t.\@\xspace}

% Evaluation shortcuts
\newcommand{\numcases}{300}
\newcommand{\numspec}{100}
\newcommand{\numeq}{100}
\newcommand{\numbug}{100}
\newcommand{\medtime}{6.5\,ms}
\newcommand{\totaltime}{1.949\,s}


% ── Packages ────────────────────────────────────────────────────────────
\usepackage{amsmath,amssymb,amsthm,mathtools}
\usepackage{tikz-cd}
\usepackage{listings}
\usepackage{xcolor}
\usepackage{xspace}
\usepackage{booktabs}
\usepackage{multirow}
\usepackage{enumitem}
\usepackage[expansion=false]{microtype}
\usepackage{hyperref}
\usepackage{cleveref}
% \usepackage{thmtools}  % disabled: not available in this TeX installation
\usepackage{float}
\usepackage{caption}
\usepackage{subcaption}
\usepackage{tabularx}
\usepackage{stmaryrd}       % \llbracket, \rrbracket
\usepackage{bbm}            % \mathbbm{1}

% ── Page geometry ───────────────────────────────────────────────────────
\usepackage[margin=1in]{geometry}

% ── Theorem environments ────────────────────────────────────────────────
\theoremstyle{plain}
\newtheorem{theorem}{Theorem}[section]
\newtheorem{lemma}[theorem]{Lemma}
\newtheorem{proposition}[theorem]{Proposition}
\newtheorem{corollary}[theorem]{Corollary}

\theoremstyle{definition}
\newtheorem{definition}[theorem]{Definition}
\newtheorem{example}[theorem]{Example}
\newtheorem{construction}[theorem]{Construction}

\theoremstyle{remark}
\newtheorem{remark}[theorem]{Remark}
\newtheorem{notation}[theorem]{Notation}

% ── Python listings style ──────────────────────────────────────────────
\definecolor{jugeoBlue}{HTML}{2563EB}
\definecolor{jugeoGreen}{HTML}{059669}
\definecolor{jugeoRed}{HTML}{DC2626}
\definecolor{jugeoPurple}{HTML}{7C3AED}
\definecolor{jugeoGray}{HTML}{6B7280}
\definecolor{jugeoBg}{HTML}{F8FAFC}

\lstdefinestyle{jugeo-python}{
  language=Python,
  basicstyle=\ttfamily\small,
  keywordstyle=\color{jugeoBlue}\bfseries,
  stringstyle=\color{jugeoGreen},
  commentstyle=\color{jugeoGray}\itshape,
  numberstyle=\tiny\color{jugeoGray},
  backgroundcolor=\color{jugeoBg},
  numbers=left,
  numbersep=6pt,
  frame=single,
  framerule=0.4pt,
  rulecolor=\color{jugeoGray!40},
  breaklines=true,
  breakatwhitespace=true,
  showstringspaces=false,
  tabsize=4,
  xleftmargin=1.5em,
  framexleftmargin=1.5em,
  morekeywords={dataclass,frozen,field,Enum,IntEnum,auto,Any,
                Mapping,Sequence,Optional,match,case},
  literate={→}{{$\to$}}1 {←}{{$\leftarrow$}}1
           {≤}{{$\leq$}}1 {≥}{{$\geq$}}1 {≠}{{$\neq$}}1
           {∈}{{$\in$}}1 {∀}{{$\forall$}}1 {∃}{{$\exists$}}1
           {⊕}{{$\oplus$}}1 {⊖}{{$\ominus$}}1,
  escapeinside={(*@}{@*)},
}
\lstset{style=jugeo-python}

\lstdefinestyle{jugeo-output}{
  basicstyle=\ttfamily\small,
  backgroundcolor=\color{jugeoBg},
  frame=single,
  framerule=0.4pt,
  rulecolor=\color{jugeoGray!40},
  breaklines=true,
  showstringspaces=false,
  xleftmargin=1.5em,
  framexleftmargin=0.5em,
  numbers=none,
}

% ── JuGeo-specific macros ──────────────────────────────────────────────

% Judgment 8-tuple
\newcommand{\J}{\mathcal{J}}
\newcommand{\Jud}[8]{(#1,\,#2,\,#3,\,#4,\,#5,\,#6,\,#7,\,#8)}
\newcommand{\JudShort}[1]{\J_{#1}}

% Components
\newcommand{\coord}{\mathsf{c}}            % coordinate
\newcommand{\prop}{\varphi}                 % proposition
\newcommand{\carrier}{\mathcal{A}}          % carrier type
\newcommand{\evid}{\mathcal{E}}             % evidence bundle
\newcommand{\oblig}{\mathcal{O}}            % obligations
\newcommand{\obstr}{\mathcal{B}}            % obstructions
\newcommand{\trust}{\mathcal{T}}            % trust annotation
\newcommand{\prov}{\Pi}                     % provenance

% Trust algebra
\newcommand{\Talg}{\mathfrak{T}}            % trust ordered algebra
\newcommand{\Eadm}{\mathcal{E}_{\mathrm{adm}}}
\newcommand{\tleq}{\preceq}                 % trust ordering
\newcommand{\tjoin}{\oplus}                 % conservative join
\newcommand{\tatten}{\ominus}               % attenuation
\newcommand{\tpromote}{\uparrow_\pi}        % promotion
\newcommand{\tdemote}{\downarrow_\chi}       % demotion

% Trust tiers
\newcommand{\Tcontra}{\ensuremath{\mathsf{CONTRADICTED}}}
\newcommand{\Tunver}{\ensuremath{\mathsf{UNVERIFIED}}}
\newcommand{\Tcopilot}{\ensuremath{\mathsf{COPILOT}}}
\newcommand{\Toracle}{\ensuremath{\mathsf{ORACLE}}}
\newcommand{\Truntime}{\ensuremath{\mathsf{RUNTIME}}}
\newcommand{\Tsolver}{\ensuremath{\mathsf{SOLVER}}}
\newcommand{\Tproof}{\ensuremath{\mathsf{PROOF}}}

% Site and geometry
\newcommand{\Site}{\mathcal{S}}             % semantic site
\newcommand{\Cov}{\mathrm{Cov}}            % covering family
\newcommand{\Ob}{\mathrm{Ob}}              % objects
\newcommand{\Mor}{\mathrm{Mor}}            % morphisms
\newcommand{\res}{\mathrm{res}}            % restriction
\newcommand{\Cech}{\check{C}}              % Čech complex
\newcommand{\Hcoh}[1]{H^{#1}}             % cohomology group

% Descent
\newcommand{\descent}{\mathrm{desc}}
\newcommand{\glue}{\mathrm{glue}}
\newcommand{\overlap}[2]{#1 \cap #2}

% Semantic moves
\newcommand{\VERIFY}{\textsc{Verify}}
\newcommand{\CONSTRUCT}{\textsc{Construct}}
\newcommand{\NORMALIZE}{\textsc{Normalize}}
\newcommand{\GLUE}{\textsc{Glue}}
\newcommand{\RESTRICT}{\textsc{Restrict}}
\newcommand{\TRANSPORT}{\textsc{Transport}}
\newcommand{\REPAIR}{\textsc{Repair}}
\newcommand{\PROMOTE}{\textsc{Promote}}
\newcommand{\CHALLENGE}{\textsc{Challenge}}
\newcommand{\ESCALATE}{\textsc{Escalate}}

% Fragments
\newcommand{\QFLIA}{\texttt{QF\_LIA}}
\newcommand{\QFLRA}{\texttt{QF\_LRA}}
\newcommand{\QFBV}{\texttt{QF\_BV}}
\newcommand{\QFUF}{\texttt{QF\_UF}}

% Misc
\newcommand{\Python}{\textsc{Python}}
\newcommand{\JuGeo}{\textsc{JuGeo}}
\newcommand{\LEAN}{\textsc{Lean}\,4}
\newcommand{\Fstar}{F$^{\star}$}
\newcommand{\Coq}{\textsc{Coq}}
\newcommand{\Dafny}{\textsc{Dafny}}

% Common shortcuts
\newcommand{\ie}{i.e.\@\xspace}
\newcommand{\eg}{e.g.\@\xspace}
\newcommand{\cf}{cf.\@\xspace}
\newcommand{\wrt}{w.r.t.\@\xspace}

% Evaluation shortcuts
\newcommand{\numcases}{300}
\newcommand{\numspec}{100}
\newcommand{\numeq}{100}
\newcommand{\numbug}{100}
\newcommand{\medtime}{6.5\,ms}
\newcommand{\totaltime}{1.949\,s}


% ── Packages ────────────────────────────────────────────────────────────
\usepackage{amsmath,amssymb,amsthm,mathtools}
\usepackage{tikz-cd}
\usepackage{listings}
\usepackage{xcolor}
\usepackage{xspace}
\usepackage{booktabs}
\usepackage{multirow}
\usepackage{enumitem}
\usepackage[expansion=false]{microtype}
\usepackage{hyperref}
\usepackage{cleveref}
% \usepackage{thmtools}  % disabled: not available in this TeX installation
\usepackage{float}
\usepackage{caption}
\usepackage{subcaption}
\usepackage{tabularx}
\usepackage{stmaryrd}       % \llbracket, \rrbracket
\usepackage{bbm}            % \mathbbm{1}

% ── Page geometry ───────────────────────────────────────────────────────
\usepackage[margin=1in]{geometry}

% ── Theorem environments ────────────────────────────────────────────────
\theoremstyle{plain}
\newtheorem{theorem}{Theorem}[section]
\newtheorem{lemma}[theorem]{Lemma}
\newtheorem{proposition}[theorem]{Proposition}
\newtheorem{corollary}[theorem]{Corollary}

\theoremstyle{definition}
\newtheorem{definition}[theorem]{Definition}
\newtheorem{example}[theorem]{Example}
\newtheorem{construction}[theorem]{Construction}

\theoremstyle{remark}
\newtheorem{remark}[theorem]{Remark}
\newtheorem{notation}[theorem]{Notation}

% ── Python listings style ──────────────────────────────────────────────
\definecolor{jugeoBlue}{HTML}{2563EB}
\definecolor{jugeoGreen}{HTML}{059669}
\definecolor{jugeoRed}{HTML}{DC2626}
\definecolor{jugeoPurple}{HTML}{7C3AED}
\definecolor{jugeoGray}{HTML}{6B7280}
\definecolor{jugeoBg}{HTML}{F8FAFC}

\lstdefinestyle{jugeo-python}{
  language=Python,
  basicstyle=\ttfamily\small,
  keywordstyle=\color{jugeoBlue}\bfseries,
  stringstyle=\color{jugeoGreen},
  commentstyle=\color{jugeoGray}\itshape,
  numberstyle=\tiny\color{jugeoGray},
  backgroundcolor=\color{jugeoBg},
  numbers=left,
  numbersep=6pt,
  frame=single,
  framerule=0.4pt,
  rulecolor=\color{jugeoGray!40},
  breaklines=true,
  breakatwhitespace=true,
  showstringspaces=false,
  tabsize=4,
  xleftmargin=1.5em,
  framexleftmargin=1.5em,
  morekeywords={dataclass,frozen,field,Enum,IntEnum,auto,Any,
                Mapping,Sequence,Optional,match,case},
  literate={→}{{$\to$}}1 {←}{{$\leftarrow$}}1
           {≤}{{$\leq$}}1 {≥}{{$\geq$}}1 {≠}{{$\neq$}}1
           {∈}{{$\in$}}1 {∀}{{$\forall$}}1 {∃}{{$\exists$}}1
           {⊕}{{$\oplus$}}1 {⊖}{{$\ominus$}}1,
  escapeinside={(*@}{@*)},
}
\lstset{style=jugeo-python}

\lstdefinestyle{jugeo-output}{
  basicstyle=\ttfamily\small,
  backgroundcolor=\color{jugeoBg},
  frame=single,
  framerule=0.4pt,
  rulecolor=\color{jugeoGray!40},
  breaklines=true,
  showstringspaces=false,
  xleftmargin=1.5em,
  framexleftmargin=0.5em,
  numbers=none,
}

% ── JuGeo-specific macros ──────────────────────────────────────────────

% Judgment 8-tuple
\newcommand{\J}{\mathcal{J}}
\newcommand{\Jud}[8]{(#1,\,#2,\,#3,\,#4,\,#5,\,#6,\,#7,\,#8)}
\newcommand{\JudShort}[1]{\J_{#1}}

% Components
\newcommand{\coord}{\mathsf{c}}            % coordinate
\newcommand{\prop}{\varphi}                 % proposition
\newcommand{\carrier}{\mathcal{A}}          % carrier type
\newcommand{\evid}{\mathcal{E}}             % evidence bundle
\newcommand{\oblig}{\mathcal{O}}            % obligations
\newcommand{\obstr}{\mathcal{B}}            % obstructions
\newcommand{\trust}{\mathcal{T}}            % trust annotation
\newcommand{\prov}{\Pi}                     % provenance

% Trust algebra
\newcommand{\Talg}{\mathfrak{T}}            % trust ordered algebra
\newcommand{\Eadm}{\mathcal{E}_{\mathrm{adm}}}
\newcommand{\tleq}{\preceq}                 % trust ordering
\newcommand{\tjoin}{\oplus}                 % conservative join
\newcommand{\tatten}{\ominus}               % attenuation
\newcommand{\tpromote}{\uparrow_\pi}        % promotion
\newcommand{\tdemote}{\downarrow_\chi}       % demotion

% Trust tiers
\newcommand{\Tcontra}{\ensuremath{\mathsf{CONTRADICTED}}}
\newcommand{\Tunver}{\ensuremath{\mathsf{UNVERIFIED}}}
\newcommand{\Tcopilot}{\ensuremath{\mathsf{COPILOT}}}
\newcommand{\Toracle}{\ensuremath{\mathsf{ORACLE}}}
\newcommand{\Truntime}{\ensuremath{\mathsf{RUNTIME}}}
\newcommand{\Tsolver}{\ensuremath{\mathsf{SOLVER}}}
\newcommand{\Tproof}{\ensuremath{\mathsf{PROOF}}}

% Site and geometry
\newcommand{\Site}{\mathcal{S}}             % semantic site
\newcommand{\Cov}{\mathrm{Cov}}            % covering family
\newcommand{\Ob}{\mathrm{Ob}}              % objects
\newcommand{\Mor}{\mathrm{Mor}}            % morphisms
\newcommand{\res}{\mathrm{res}}            % restriction
\newcommand{\Cech}{\check{C}}              % Čech complex
\newcommand{\Hcoh}[1]{H^{#1}}             % cohomology group

% Descent
\newcommand{\descent}{\mathrm{desc}}
\newcommand{\glue}{\mathrm{glue}}
\newcommand{\overlap}[2]{#1 \cap #2}

% Semantic moves
\newcommand{\VERIFY}{\textsc{Verify}}
\newcommand{\CONSTRUCT}{\textsc{Construct}}
\newcommand{\NORMALIZE}{\textsc{Normalize}}
\newcommand{\GLUE}{\textsc{Glue}}
\newcommand{\RESTRICT}{\textsc{Restrict}}
\newcommand{\TRANSPORT}{\textsc{Transport}}
\newcommand{\REPAIR}{\textsc{Repair}}
\newcommand{\PROMOTE}{\textsc{Promote}}
\newcommand{\CHALLENGE}{\textsc{Challenge}}
\newcommand{\ESCALATE}{\textsc{Escalate}}

% Fragments
\newcommand{\QFLIA}{\texttt{QF\_LIA}}
\newcommand{\QFLRA}{\texttt{QF\_LRA}}
\newcommand{\QFBV}{\texttt{QF\_BV}}
\newcommand{\QFUF}{\texttt{QF\_UF}}

% Misc
\newcommand{\Python}{\textsc{Python}}
\newcommand{\JuGeo}{\textsc{JuGeo}}
\newcommand{\LEAN}{\textsc{Lean}\,4}
\newcommand{\Fstar}{F$^{\star}$}
\newcommand{\Coq}{\textsc{Coq}}
\newcommand{\Dafny}{\textsc{Dafny}}

% Common shortcuts
\newcommand{\ie}{i.e.\@\xspace}
\newcommand{\eg}{e.g.\@\xspace}
\newcommand{\cf}{cf.\@\xspace}
\newcommand{\wrt}{w.r.t.\@\xspace}

% Evaluation shortcuts
\newcommand{\numcases}{300}
\newcommand{\numspec}{100}
\newcommand{\numeq}{100}
\newcommand{\numbug}{100}
\newcommand{\medtime}{6.5\,ms}
\newcommand{\totaltime}{1.949\,s}

% experiment-data.tex — AUTO-GENERATED from real jugeo CLI runs
% DO NOT EDIT — regenerate with: python3 experiments/run_universal_metrics.py
% Generated from 30 programs across 6 domains

\newcommand{\expTotalPrograms}{30}
\newcommand{\expVerified}{30}
\newcommand{\expAccuracy}{100.0\%}
\newcommand{\expCoordsMin}{2}
\newcommand{\expCoordsMax}{10}
\newcommand{\expCoordsMean}{5.2}
\newcommand{\expCoordsMedian}{5.5}
\newcommand{\expPropsMin}{4}
\newcommand{\expPropsMax}{20}
\newcommand{\expPropsMean}{10.4}
\newcommand{\expPropsSum}{311}
\newcommand{\expPropsOkSum}{310}
\newcommand{\expTimeMin}{3.506\,s}
\newcommand{\expTimeMax}{6.714\,s}
\newcommand{\expTimeMean}{4.64\,s}
\newcommand{\expTimeMedian}{4.508\,s}
\newcommand{\expTimeTotal}{139.204\,s}
\newcommand{\expObstructionTotal}{0}
\newcommand{\expHOneZero}{30}
\newcommand{\expDomainCount}{6}
\newcommand{\expTrustCOPILOTSUGGESTED}{30}
\newcommand{\expDomainDsCount}{5}
\newcommand{\expDomainDsVerified}{5}
\newcommand{\expDomainDsAvgCoords}{7.6}
\newcommand{\expDomainDsAvgProps}{15.2}
\newcommand{\expDomainMathCount}{5}
\newcommand{\expDomainMathVerified}{5}
\newcommand{\expDomainMathAvgCoords}{4.8}
\newcommand{\expDomainMathAvgProps}{9.4}
\newcommand{\expDomainSmCount}{5}
\newcommand{\expDomainSmVerified}{5}
\newcommand{\expDomainSmAvgCoords}{7.6}
\newcommand{\expDomainSmAvgProps}{15.2}
\newcommand{\expDomainSortCount}{5}
\newcommand{\expDomainSortVerified}{5}
\newcommand{\expDomainSortAvgCoords}{2.4}
\newcommand{\expDomainSortAvgProps}{4.8}
\newcommand{\expDomainStrCount}{5}
\newcommand{\expDomainStrVerified}{5}
\newcommand{\expDomainStrAvgCoords}{3.2}
\newcommand{\expDomainStrAvgProps}{6.4}
\newcommand{\expDomainWebCount}{5}
\newcommand{\expDomainWebVerified}{5}
\newcommand{\expDomainWebAvgCoords}{5.6}
\newcommand{\expDomainWebAvgProps}{11.2}

% subsystem-data.tex — AUTO-GENERATED from real jugeo subsystem experiments
% DO NOT EDIT — regenerate with: python3 experiments/run_subsystem_experiments.py

\newcommand{\subCliTotal}{10}
\newcommand{\subCliVerified}{9}
\newcommand{\subCliAccuracy}{90.0\%}
\newcommand{\subCliMeanTime}{4.132\,s}
\newcommand{\subCliMinTime}{3.472\,s}
\newcommand{\subCliMaxTime}{5.372\,s}
\newcommand{\subCliTotalCoords}{54}
\newcommand{\subCliTotalProps}{107}
\newcommand{\subCliTotalPropsOk}{106}
\newcommand{\subSiteCalls}{90}
\newcommand{\subSiteSuccessful}{69}
\newcommand{\subSiteSuccessRate}{76.7\%}
\newcommand{\subSiteMeanTime}{0.0001\,s}
\newcommand{\subSiteTotalTime}{0.005\,s}
\newcommand{\subSpecificationsatisfactionSmallTime}{0.0003\,s}
\newcommand{\subSpecificationsatisfactionMediumTime}{0.0001\,s}
\newcommand{\subSpecificationsatisfactionLargeTime}{0.0001\,s}
\newcommand{\subInhabitantfleetSmallTime}{0.0\,s}
\newcommand{\subInhabitantfleetMediumTime}{0.0\,s}
\newcommand{\subInhabitantfleetLargeTime}{0.0\,s}
\newcommand{\subTheoremecologySmallTime}{0.0\,s}
\newcommand{\subTheoremecologyMediumTime}{0.0\,s}
\newcommand{\subTheoremecologyLargeTime}{0.0\,s}
\newcommand{\subStatespaceexplorationSmallTime}{0.0\,s}
\newcommand{\subStatespaceexplorationMediumTime}{0.0\,s}
\newcommand{\subStatespaceexplorationLargeTime}{0.0\,s}
\newcommand{\subRepairsemanticsSmallTime}{0.0\,s}
\newcommand{\subRepairsemanticsMediumTime}{0.0\,s}
\newcommand{\subRepairsemanticsLargeTime}{0.0\,s}
\newcommand{\subReplaygluingSmallTime}{0.0\,s}
\newcommand{\subReplaygluingMediumTime}{0.0\,s}
\newcommand{\subReplaygluingLargeTime}{0.0\,s}
\newcommand{\subSemanticfuturesSmallTime}{0.0\,s}
\newcommand{\subSemanticfuturesMediumTime}{0.0\,s}
\newcommand{\subSemanticfuturesLargeTime}{0.0\,s}
\newcommand{\subHypercovertreatySmallTime}{0.0\,s}
\newcommand{\subHypercovertreatyMediumTime}{0.0\,s}
\newcommand{\subHypercovertreatyLargeTime}{0.0\,s}
\newcommand{\subEncodeforsolverSmallTime}{0.0026\,s}
\newcommand{\subEncodeforsolverMediumTime}{0.0002\,s}
\newcommand{\subEncodeforsolverLargeTime}{0.0001\,s}
\newcommand{\subInterfaceroutingSmallTime}{0.0\,s}
\newcommand{\subInterfaceroutingMediumTime}{0.0\,s}
\newcommand{\subInterfaceroutingLargeTime}{0.0\,s}
\newcommand{\subOrchestrateverificationSmallTime}{0.0002\,s}
\newcommand{\subOrchestrateverificationMediumTime}{0.0\,s}
\newcommand{\subOrchestrateverificationLargeTime}{0.0\,s}
\newcommand{\subRunfulldescentSmallTime}{0.0\,s}
\newcommand{\subRunfulldescentMediumTime}{0.0\,s}
\newcommand{\subRunfulldescentLargeTime}{0.0\,s}
\newcommand{\subKernellifecycleSmallTime}{0.0\,s}
\newcommand{\subKernellifecycleMediumTime}{0.0\,s}
\newcommand{\subKernellifecycleLargeTime}{0.0\,s}
\newcommand{\subDiscoverypipelineSmallTime}{0.0\,s}
\newcommand{\subDiscoverypipelineMediumTime}{0.0\,s}
\newcommand{\subDiscoverypipelineLargeTime}{0.0\,s}
\newcommand{\subAnalogytransportSmallTime}{0.0\,s}
\newcommand{\subAnalogytransportMediumTime}{0.0\,s}
\newcommand{\subAnalogytransportLargeTime}{0.0\,s}
\newcommand{\subFormalcoresiteSmallTime}{0.0001\,s}
\newcommand{\subFormalcoresiteMediumTime}{0.0\,s}
\newcommand{\subFormalcoresiteLargeTime}{0.0\,s}
\newcommand{\subProblematlasSmallTime}{0.0\,s}
\newcommand{\subProblematlasMediumTime}{0.0\,s}
\newcommand{\subProblematlasLargeTime}{0.0\,s}
\newcommand{\subPublicalignmentSmallTime}{0.0\,s}
\newcommand{\subPublicalignmentMediumTime}{0.0\,s}
\newcommand{\subPublicalignmentLargeTime}{0.0\,s}
\newcommand{\subRegimebootstrappingSmallTime}{0.0\,s}
\newcommand{\subRegimebootstrappingMediumTime}{0.0\,s}
\newcommand{\subRegimebootstrappingLargeTime}{0.0\,s}
\newcommand{\subRelationalrefinementSmallTime}{0.0\,s}
\newcommand{\subRelationalrefinementMediumTime}{0.0\,s}
\newcommand{\subRelationalrefinementLargeTime}{0.0\,s}
\newcommand{\subSemanticclosureSmallTime}{0.0\,s}
\newcommand{\subSemanticclosureMediumTime}{0.0\,s}
\newcommand{\subSemanticclosureLargeTime}{0.0\,s}
\newcommand{\subTheoremeconomicsSmallTime}{0.0001\,s}
\newcommand{\subTheoremeconomicsMediumTime}{0.0\,s}
\newcommand{\subTheoremeconomicsLargeTime}{0.0\,s}
\newcommand{\subBenchmarksuiteSmallTime}{0.0\,s}
\newcommand{\subBenchmarksuiteMediumTime}{0.0\,s}
\newcommand{\subBenchmarksuiteLargeTime}{0.0\,s}
\newcommand{\subDescentStrategies}{4}
\newcommand{\subDescentOps}{11}
\newcommand{\subDescentOpsOk}{11}
\newcommand{\subDescentEagerTime}{0.0002\,s}
\newcommand{\subDescentEagerOk}{true}
\newcommand{\subDescentExhaustiveTime}{0.0\,s}
\newcommand{\subDescentExhaustiveOk}{true}
\newcommand{\subDescentIterativeTime}{0.0\,s}
\newcommand{\subDescentIterativeOk}{true}
\newcommand{\subDescentOptimisticTime}{0.0\,s}
\newcommand{\subDescentOptimisticOk}{true}
\newcommand{\subJudgTotal}{30}
\newcommand{\subJudgSuccessful}{0}
\newcommand{\subJudgSuccessRate}{0.0\%}
\newcommand{\subJudgMeanTimeUs}{7.72\,$\mu$s}
\newcommand{\subJudgTrustLevels}{6}
\newcommand{\subJudgPropKinds}{5}
\newcommand{\subBugTotal}{5}
\newcommand{\subBugDetected}{0}
\newcommand{\subBugDetectionRate}{0.0\%}
\newcommand{\subBugFalsePositiveRate}{0.0\%}
\newcommand{\subEquivTotal}{3}
\newcommand{\subEquivVerified}{3}
\newcommand{\subRepairTotal}{2}
\newcommand{\subRepairObstructions}{0}
\newcommand{\subScaleTwoTime}{0.0001\,s}
\newcommand{\subScaleFourTime}{0.0\,s}
\newcommand{\subScaleEightTime}{0.0\,s}
\newcommand{\subScaleTwelveTime}{0.0\,s}
\newcommand{\subScaleSixteenTime}{0.0\,s}
\newcommand{\subScaleTwentyTime}{0.0\,s}
\newcommand{\subTotalExperimentTime}{95.6\,s}

% data-paper27.tex — AUTO-GENERATED by exp27_certificate_chains.py
% DO NOT EDIT — regenerate with: python3 experiments/exp27_certificate_chains.py

\newcommand{\ppTwentysevenSpecCount}{4}
\newcommand{\ppTwentysevenSpecMean}{29749.07\,\text{ms}}
\newcommand{\ppTwentysevenSpecPnn}{29158.79\,\text{ms}}
\newcommand{\ppTwentysevenSpecMax}{38435.38\,\text{ms}}
\newcommand{\ppTwentysevenEquivCount}{4}
\newcommand{\ppTwentysevenEquivMean}{24565.53\,\text{ms}}
\newcommand{\ppTwentysevenEquivPnn}{24594.32\,\text{ms}}
\newcommand{\ppTwentysevenEquivMax}{26518.73\,\text{ms}}
\newcommand{\ppTwentysevenBugDetCount}{4}
\newcommand{\ppTwentysevenBugDetMean}{22072.72\,\text{ms}}
\newcommand{\ppTwentysevenBugDetPnn}{21732.74\,\text{ms}}
\newcommand{\ppTwentysevenBugDetMax}{24949.48\,\text{ms}}
\newcommand{\ppTwentysevenAllCount}{12}
\newcommand{\ppTwentysevenAllMean}{25462.44\,\text{ms}}
\newcommand{\ppTwentysevenAllPnn}{29158.79\,\text{ms}}
\newcommand{\ppTwentysevenAllMax}{38435.38\,\text{ms}}
\newcommand{\ppTwentysevenTotalPrograms}{12}
\newcommand{\ppTwentysevenTotalCertificates}{47}
\newcommand{\ppTwentysevenTotalMorphisms}{50}
\newcommand{\ppTwentysevenVerifiedCount}{12}
\newcommand{\ppTwentysevenVerifiedRate}{100.0\%}
\newcommand{\ppTwentysevenTotalSections}{47}
\newcommand{\ppTwentysevenTotalOverlaps}{83}
\newcommand{\ppTwentysevenAvgChainLen}{3.9}


\title{Certificate Chains: Compositional Proof-Carrying Code Architecture}

\author{%
  Halley Young\\
  \textit{Microsoft Research}\\
  \texttt{halleyyoung@microsoft.com}
}
\date{\today}

% ── Paper-local macros ──────────────────────────────────────────────
\newcommand{\Cert}{\mathsf{Cert}}              % certificate
\newcommand{\Chain}{\mathsf{Chain}}            % certificate chain
\newcommand{\CA}{\mathsf{CA}}                  % certificate authority
\newcommand{\CV}{\mathsf{CV}}                  % certificate verifier
\newcommand{\CM}{\mathsf{CM}}                  % certificate merger
\newcommand{\CP}{\mathsf{CP}}                  % certificate projection
\newcommand{\CS}{\mathsf{CS}}                  % certificate serializer
\newcommand{\certid}{\mathit{id}}              % certificate identifier
\newcommand{\certcoord}{\mathit{coord}}        % certificate coordinate
\newcommand{\certprops}{\mathit{props}}        % verified propositions
\newcommand{\certtrust}{\mathit{tl}}           % trust level
\newcommand{\certissuer}{\mathit{iss}}         % issuer
\newcommand{\certsig}{\mathit{sig}}            % signature hash
\newcommand{\certresid}{\mathit{res}}          % residual obligations
\newcommand{\certobstr}{\mathit{obs}}          % obstructions
\newcommand{\certexpiry}{\mathit{exp}}         % expiry
\newcommand{\Valid}{\mathrm{Valid}}            % validity predicate
\newcommand{\Issued}{\mathrm{Issued}}          % issued registry
\newcommand{\Revoked}{\mathrm{Revoked}}        % revoked set
\newcommand{\TL}{\mathbf{TL}}                  % trust level set
\newcommand{\tlmin}{\wedge_{\TL}}             % trust minimum
\newcommand{\Sigalg}{\mathcal{H}}              % signature algorithm (SHA-256)
\newcommand{\Prop}{\mathsf{Prop}}              % proposition type
\newcommand{\AuthSet}{\mathcal{A}\!\mathit{uth}} % authority set
\newcommand{\nss}{\mathrm{NSS}}                % no-silent-strengthening
\newcommand{\cover}{\mathrm{cover}}            % coverage predicate
\renewcommand{\merge}{\mathrm{merge}}            % merge operation
\newcommand{\proj}{\mathrm{proj}}              % projection
\newcommand{\serial}{\mathrm{serial}}          % serialization
\newcommand{\Pending}{\mathsf{PENDING}}
\newcommand{\Settled}{\mathsf{SETTLED}}
\newcommand{\Obstructed}{\mathsf{OBSTRUCTED}}
\newcommand{\Revk}{\mathsf{REVOKED}}
\newcommand{\Expired}{\mathsf{EXPIRED}}
\newcommand{\StatusSet}{\{\Pending,\Settled,\Obstructed,\Revk,\Expired\}}
\newcommand{\etc}{etc.\@\xspace}

% ── BibTeX-less references ──────────────────────────────────────────
\newcommand{\citeNec}{[NL97]}
\newcommand{\citeApp}{[App01]}
\newcommand{\citeFoundPCC}{[Nec97]}
\newcommand{\citeTAL}{[MCG99]}
\newcommand{\citeProofWeb}{[VPL05]}
\newcommand{\citeLCC}{[Mil07]}
\newcommand{\citeCoqExt}{[BC04]}
\newcommand{\citeXcert}{[RFC5280]}
\newcommand{\citePKI}{[Ad01]}
\newcommand{\citeHAT}{[HAT04]}

% ───────────────────────────────────────────────────────────────────
\begin{document}
\maketitle

% =====================================================================
\begin{abstract}
We present a \emph{certificate chain} architecture for compositional
proof-carrying code (PCC) in the \JuGeo{} system.  The paper develops a
chain model for aggregating local verification results, weakest-link
trust semantics, and chain-level diagnostics.  The repository-backed
evaluation covers 12 benchmark programs, 47 certificates, and 83
overlaps.  On that benchmark, the mean end-to-end processing time is
25462.44\,ms, with per-suite means between roughly 22\,s and 30\,s.
Accordingly, this audited version makes timing claims only at those
measured scales and does not retain unsupported claims about
sub-millisecond verification or unmeasured forgery experiments.
\end{abstract}

\newpage

% =====================================================================
\section{Introduction}\label{sec:intro}
% =====================================================================

Proof-carrying code (PCC), introduced by Necula and Lee
\citeNec\,\citeFoundPCC, attaches a machine-checkable proof to a binary,
allowing a receiver to verify safety properties without trusting the
producer.
Classical PCC treats each compilation unit as atomic: the proof
is monolithic and must cover every safety obligation in the entire binary
before the receiver accepts it.
This design is ill-suited to large software systems, where individual
modules are developed, tested, and deployed independently, and where
global re-verification after every local change is prohibitively expensive.

The \JuGeo{} framework addresses this scalability gap through its
\emph{judgment geometry}: each coordinate in the codebase carries an
8-tuple judgment $(c,\varphi,\carrier,\evid,\oblig,\obstr,\trust,\prov)$
whose evidence bundle $\evid$ and trust annotation $\trust$ can be
checked locally.
What has been missing is a \emph{certificate layer}: a protocol by which
local verifications are recorded, signed, chained together, and shipped
alongside code in a form that downstream consumers can efficiently verify.

This paper presents that layer.
Our main contributions are:

\begin{enumerate}[leftmargin=2em]
  \item \textbf{A certificate data model} (\S\ref{sec:cert}): a frozen,
        content-addressed record carrying a coordinate, verified propositions,
        trust level, residual obligations, obstructions, issuer identity, and
        a SHA-256 signature hash.

  \item \textbf{Chain construction and the weakest-link trust semantics}
        (\S\ref{sec:chain}): \texttt{CertificateChain} assembles an ordered
        sequence of certificates; the chain's aggregate trust is the minimum
        trust of any link, ensuring conservative global guarantees.

  \item \textbf{A delegated authority model} (\S\ref{sec:authority}):
        \texttt{CertificateAuthority} governs issuance, revocation, and renewal
        within a configurable delegation depth, preventing unauthorized parties
        from injecting certificates into a chain.

  \item \textbf{A chain-verification algorithm with the no-silent-strengthening
        invariant} (\S\ref{sec:verify}): \texttt{CertificateVerifier} checks
        every link for structural consistency, trust monotonicity, and the
        absence of inflated claims.

  \item \textbf{Merger and projection} (\S\ref{sec:merge}): \texttt{CertificateMerger}
        reconciles overlapping certificates under three strategies;
        \texttt{CertificateProjection} produces role-appropriate views while
        preserving faithfulness.

  \item \textbf{Serialization} (\S\ref{sec:serial}):
        \texttt{CertificateSerializer} provides full round-trip JSON
        fidelity for shipping certificates with compiled artifacts.

  \item \textbf{Chain Integrity theorem} (\S\ref{sec:theorem}):
        formally proved in \LEAN{}, establishing soundness and
        anti-forgery guarantees.
\end{enumerate}

\paragraph{Related work.}
Typed assembly language (TAL) \citeTAL{} and proof-carrying
authentication \citeHAT{} share our goal of shipping machine-checkable
guarantees with code; unlike those systems, \JuGeo{} operates at the
level of rich judgment geometry rather than type systems alone.
X.509 certificate chains \citeXcert{} inspire our authority model, but
our chains attest to semantic verification properties, not merely identity.
Foundational PCC \citeApp{} and LCC \citeLCC{} pursue similar compositionality
goals; our contribution is a \Python{}-native, geometry-aware realization
integrated with the full \JuGeo{} evidence pipeline.

% =====================================================================
\section{Certificate Structure}\label{sec:cert}
% =====================================================================

\subsection{The \texttt{Certificate} Record}

A \emph{certificate} is a signed, content-addressed attestation that a
specific set of propositions has been verified at a given coordinate.

\begin{definition}[Certificate]\label{def:cert}
A \emph{certificate} is an immutable tuple
\[
  \Cert = (\certid,\,\certcoord,\,\certprops,\,\certtrust,\,
           \certissuer,\,\certsig,\,\certresid,\,\certobstr,\,\certexpiry)
\]
where:
\begin{itemize}
  \item $\certid \in \mathsf{UUID}$: a unique identifier.
  \item $\certcoord \in \mathsf{Coordinate}$: the codebase location being attested.
  \item $\certprops \subseteq \Prop$: the finite set of verified propositions.
  \item $\certtrust \in \TL$: the trust level ($\Pending \le \Settled \le \cdots$).
  \item $\certissuer \in \AuthSet$: the issuing authority's name.
  \item $\certsig = \Sigalg(\certcoord \| \certprops \| \certtrust \| \certissuer)$:
        a SHA-256 content fingerprint.
  \item $\certresid$: a (possibly empty) tuple of residual obligations not yet
        discharged by $\certprops$.
  \item $\certobstr$: a (possibly empty) tuple of blocking conditions.
  \item $\certexpiry \in \mathsf{Timestamp} \cup \{\bot\}$: optional expiry.
\end{itemize}
\end{definition}

Immutability is enforced by Python's \texttt{@dataclass(frozen=True)};
the signature hash~$\certsig$ is computed automatically by
\texttt{CertificateBuilder.sign()} before \texttt{build()} is called,
binding the content to the issuer identity.

\subsection{Certificate Status}\label{sec:status}

Every certificate carries a lifecycle status drawn from
$\StatusSet$, determined at query time by
\texttt{CertificateVerifier.verify()}:

\begin{center}
\begin{tabular}{lp{9cm}}
\toprule
\textbf{Status} & \textbf{Semantics} \\
\midrule
$\Pending$    & Issued but not yet independently verified. \\
$\Settled$    & Passes all \texttt{CertificateVerifier} checks; propositions confirmed. \\
$\Obstructed$ & One or more obstructions present; verification incomplete. \\
$\Revk$       & Authority has revoked the certificate; must not be trusted. \\
$\Expired$    & Current time exceeds $\certexpiry$; no longer valid. \\
\bottomrule
\end{tabular}
\end{center}

\begin{definition}[Validity]\label{def:valid}
A certificate $c$ is \emph{valid} if its status is $\Settled$:
\[
  \Valid(c) \;\iff\; c.\certobstr = \emptyset \;\wedge\; c \notin \Revoked
             \;\wedge\; \bigl(c.\certexpiry = \bot \;\vee\; \mathsf{now} < c.\certexpiry\bigr).
\]
\end{definition}

\subsection{\texttt{CertificateBuilder}: Fluent Construction}

The builder pattern enforces that all mandatory fields are populated
before a certificate is created.
The chain of calls
\begin{lstlisting}[style=jugeo-python]
cert = (CertificateBuilder()
        .for_coordinate(coord)
        .add_verified(prop)
        .set_trust(TrustLevel.VERIFIED)
        .set_issuer(authority_name)
        .set_evidence_summary(summary)
        .sign()
        .build())
\end{lstlisting}
ensures that $\certsig$ is always consistent with the content fields;
calling \texttt{build()} without \texttt{sign()} raises
\texttt{CertificateBuildError}.

\subsection{Trust Levels}

\JuGeo{} uses a three-tier trust lattice for certificates (distinct from
the eight-tier judgment trust in \texttt{jugeo-common}):

\begin{center}
\begin{tabular}{lll}
\toprule
\textbf{Level} & \textbf{Name} & \textbf{Meaning} \\
\midrule
1 & \texttt{PROPOSAL}  & Draft; not yet reviewed. \\
2 & \texttt{REVIEWED}  & Human or tool review completed. \\
3 & \texttt{VERIFIED}  & Formally or mechanically verified. \\
\bottomrule
\end{tabular}
\end{center}

The ordering $1 \le 2 \le 3$ is total; the conservative join is $\min$.

% =====================================================================
\section{Chain Construction}\label{sec:chain}
% =====================================================================

\subsection{\texttt{CertificateChain}}

A \emph{certificate chain} is an ordered list of certificates
$[c_1, c_2, \ldots, c_n]$ where each $c_i$ attests to a local
verification at coordinate $\certcoord_i$.  The chain as a whole
constitutes a global verification of $\bigcup_i \certprops_i$ across
$\{\certcoord_1, \ldots, \certcoord_n\}$.

\begin{definition}[Certificate Chain]\label{def:chain}
A \emph{certificate chain} over a set of coordinates
$\{\certcoord_1, \ldots, \certcoord_n\}$ is an ordered sequence
$\Chain = [c_1, \ldots, c_n]$ of certificates satisfying:
\begin{enumerate}
  \item \emph{Coverage}: $c_i.\certcoord = \certcoord_i$ for each $i$.
  \item \emph{Integrity}: each $c_i$ carries a valid signature hash.
  \item \emph{Trust floor}: $\mathrm{tf}(\Chain) = \min_{i} c_i.\certtrust$.
\end{enumerate}
\end{definition}

\subsection{Weakest-Link Trust Semantics}

The \emph{trust floor} of a chain is conservative:
\[
  \mathrm{tf}([c_1,\ldots,c_n]) = c_1.\certtrust \;\tlmin\; c_2.\certtrust
    \;\tlmin\; \cdots \;\tlmin\; c_n.\certtrust.
\]
This mirrors the behavior of the trust algebra in \texttt{jugeo-common}:
joining evidence over a cover is conservative.
A chain whose weakest link is a \texttt{PROPOSAL}-level certificate
cannot be promoted to \texttt{VERIFIED} by the presence of stronger links.

\begin{proposition}[Weakest-Link Monotonicity]\label{prop:weakest-link}
Let $\Chain = [c_1,\ldots,c_n]$ and $\Chain' = [c_1,\ldots,c_n,c_{n+1}]$.
Then $\mathrm{tf}(\Chain') \le \mathrm{tf}(\Chain)$.
\end{proposition}
\begin{proof}
$\mathrm{tf}(\Chain') = \mathrm{tf}(\Chain) \;\tlmin\; c_{n+1}.\certtrust \le \mathrm{tf}(\Chain)$
by the definition of $\tlmin$.
\end{proof}

\subsection{Chain Operations}

\begin{center}
\begin{tabular}{lp{8.5cm}}
\toprule
\textbf{Method} & \textbf{Semantics} \\
\midrule
\texttt{verify\_chain()} & Returns \texttt{True} iff every link is valid and
  the signature hashes are consistent. \\
\texttt{trust\_floor()} & Returns $\mathrm{tf}(\Chain)$. \\
\texttt{weakest\_link()} & Returns $\arg\min_i c_i.\certtrust$. \\
\texttt{is\_complete()} & Checks that every required coordinate has a link. \\
\texttt{gaps()} & Returns coordinates with no covering certificate. \\
\texttt{extend\_chain(c)} & Appends $c$ to the chain (immutable: returns new chain). \\
\texttt{coordinates()} & Returns $\{\certcoord_1,\ldots,\certcoord_n\}$. \\
\texttt{total\_residuals()} & Returns $\sum_i |\certresid_i|$. \\
\texttt{total\_obstructions()} & Returns $\sum_i |\certobstr_i|$. \\
\bottomrule
\end{tabular}
\end{center}

\subsection{Compositionality}

The key compositionality property of chains is that the global
verification they establish can be checked \emph{link by link}:
no link needs information from any other link beyond what is
encoded in its own certificate.
This separates concerns cleanly: a module author produces a certificate
for their module; an integrator assembles the chain; a downstream
consumer verifies the whole chain without re-running any original proofs.

% =====================================================================
\section{Authority Model}\label{sec:authority}
% =====================================================================

\subsection{\texttt{CertificateAuthority}}

A \texttt{CertificateAuthority} (CA) is the entity authorized to issue,
revoke, and renew certificates.

\begin{definition}[Certificate Authority]\label{def:ca}
A \emph{certificate authority} is a triple
$\CA = (\mathit{name}, \AuthSet_{\mathrm{trusted}}, d_{\max})$
where:
\begin{itemize}
  \item $\mathit{name} \in \AuthSet$: the authority's identity.
  \item $\AuthSet_{\mathrm{trusted}} \subseteq \AuthSet$: the set of
        issuers whose certificates this authority will co-sign or delegate to.
  \item $d_{\max} \in \mathbb{N}$ (default 5): maximum delegation depth.
\end{itemize}
The authority maintains:
\begin{itemize}
  \item $\Issued$: a registry mapping certificate IDs to certificates.
  \item $\Revoked$: the set of revoked certificate IDs.
\end{itemize}
\end{definition}

\subsection{Issuance and Revocation}

\texttt{CA.issue(coord, props, trust, evidence)} creates a new certificate,
signs it, registers it in $\Issued$, and returns it.
The implementation enforces:

\begin{enumerate}
  \item \emph{Identity binding}: the issuer field of every issued certificate
        equals $\CA.\mathit{name}$.
  \item \emph{No re-issuance of revoked certificates}: if a coordinate already
        has a revoked certificate, a new issuance does not restore the old ID.
  \item \emph{Delegation depth}: the CA will not issue certificates on behalf
        of an authority whose delegation depth exceeds $d_{\max}$.
\end{enumerate}

Revocation is irrevocable: once $\certid \in \Revoked$, the certificate
is invalid regardless of its content.
Renewal creates a fresh certificate with a new $\certid$ and refreshed
$\certexpiry$; it does not modify the revocation set.

\subsection{Trust Boundary}

Only certificates whose $\certissuer \in \AuthSet_{\mathrm{trusted}}$ are
accepted by a CA during chain validation.
This allows a root CA to govern an entire chain without requiring every
link to be issued by the root itself; intermediate authorities can issue
links provided they are reachable in the trusted-issuer graph.

\begin{proposition}[Authority Closure]\label{prop:auth-closure}
If $\CA$ trusts $\CA'$ and $\CA'$ trusts $\CA''$, then every certificate
issued by $\CA''$ within delegation depth $d_{\max}$ is accepted by
\texttt{CertificateVerifier} when validating chains for $\CA$.
\end{proposition}

% =====================================================================
\section{Verification}\label{sec:verify}
% =====================================================================

\subsection{\texttt{CertificateVerifier}}

The verifier is an independent component that accepts or rejects a
certificate or chain without consulting the issuing authority at runtime.
Its design enforces the \emph{no-silent-strengthening} (NSS) invariant
from the \JuGeo{} trust algebra: a certificate may not claim more than
its evidence supports.

\subsection{Per-Link Checks}

For each link $c_i$ in the chain, \texttt{CertificateVerifier} performs:

\begin{enumerate}[leftmargin=2em,label=(\roman*)]
  \item \textbf{Signature check}: recompute $\Sigalg(\certcoord_i \| \certprops_i \| \certtrust_i \| \certissuer_i)$ and compare with $c_i.\certsig$.
  \item \textbf{Expiry check}: $\mathsf{now} < c_i.\certexpiry$ (or $c_i.\certexpiry = \bot$).
  \item \textbf{Revocation check}: $c_i.\certid \notin \Revoked$.
  \item \textbf{Evidence consistency}: the evidence summary is non-empty and
        references at least one of the verified propositions.
  \item \textbf{Trust consistency}: $c_i.\certtrust$ is not higher than the
        trust level justified by $c_i$'s evidence kind (solver evidence cannot
        claim $\Tproof$-level trust, \etc).
  \item \textbf{No-silent-strengthening}: $|\certresid_i|$ honestly accounts for
        every unresolved obligation; the verifier computes an independent
        obligation count from the evidence and rejects links where the
        claimed residual count is lower.
  \item \textbf{Obstruction honesty}: obstructions are not silently omitted.
\end{enumerate}

\subsection{Chain-Level Checks}

After per-link checks, the verifier performs:
\begin{enumerate}[leftmargin=2em,label=(\roman*)]
  \item \textbf{Coverage completeness}: the union $\bigcup_i \certcoord_i$
        covers the required coordinate set (provided by the caller).
  \item \textbf{Trust floor}: $\mathrm{tf}(\Chain) \ge \certtrust_{\min}$
        required by the consuming context.
  \item \textbf{Global residuals}: $\sum_i |\certresid_i| = 0$ for chains
        claiming complete verification.
\end{enumerate}

\subsection{Copilot-Assisted Verification}

\texttt{CertificateVerifier.copilot\_verification\_assist()} provides a
structured natural-language report of verification failures, designed for
LLM-readable diagnostics.  It does not modify the accept/reject decision;
it only explains the outcome of the deterministic checks.

\subsection{Chain Validation Algorithm}\label{sec:chain-val-alg}

We present the full chain-validation algorithm as pseudocode, making
explicit the ordering of checks and early-termination conditions.

\begin{definition}[Chain Validation Procedure]\label{def:chain-val}
Given a chain $\Chain = [c_1, \ldots, c_n]$, a certificate authority $\CA$,
a required coordinate set $S$, and a minimum trust threshold
$\certtrust_{\min}$, the validation procedure
$\textsc{ValidateChain}(\Chain, \CA, S, \certtrust_{\min})$
returns either $(\mathsf{True}, \mathit{report})$ or
$(\mathsf{False}, \mathit{failures})$.
\end{definition}

\begin{center}
\begin{tabular}{l}
\toprule
\textbf{Algorithm}: $\textsc{ValidateChain}(\Chain, \CA, S, \certtrust_{\min})$ \\
\midrule
\textbf{Input}: Chain $\Chain = [c_1, \ldots, c_n]$, authority $\CA$, \\
\quad required coordinates $S$, minimum trust $\certtrust_{\min}$ \\
\textbf{Output}: $(\mathsf{True}, \mathit{report})$ or $(\mathsf{False}, \mathit{failures})$ \\
\midrule
1.\quad $\mathit{failures} \gets []$ \\
2.\quad \textbf{for} $i \gets 1$ \textbf{to} $n$ \textbf{do} \\
3.\quad\quad $h \gets \Sigalg(c_i.\certcoord \| c_i.\certprops \| c_i.\certtrust \| c_i.\certissuer)$ \\
4.\quad\quad \textbf{if} $h \ne c_i.\certsig$ \textbf{then} append $\mathtt{sig\_mismatch}(i)$ to $\mathit{failures}$ \\
5.\quad\quad \textbf{if} $c_i.\certexpiry \ne \bot$ \textbf{and} $\mathsf{now} \ge c_i.\certexpiry$ \\
6.\quad\quad\quad \textbf{then} append $\mathtt{expired}(i)$ to $\mathit{failures}$ \\
7.\quad\quad \textbf{if} $c_i.\certid \in \CA.\Revoked$ \\
8.\quad\quad\quad \textbf{then} append $\mathtt{revoked}(i)$ to $\mathit{failures}$ \\
9.\quad\quad \textbf{if} $c_i.\certissuer \notin \CA.\AuthSet_{\mathrm{trusted}}$ \\
10.\quad\quad\quad \textbf{then} append $\mathtt{untrusted\_issuer}(i)$ to $\mathit{failures}$ \\
11.\quad\quad \textbf{if} $|\certresid_i| < \textsc{IndepResidCount}(c_i)$ \\
12.\quad\quad\quad \textbf{then} append $\mathtt{nss\_violation}(i)$ to $\mathit{failures}$ \\
13.\quad \textbf{end for} \\
14.\quad $C \gets \bigcup_{i=1}^{n} \{c_i.\certcoord\}$ \\
15.\quad \textbf{if} $S \not\subseteq C$ \\
16.\quad\quad \textbf{then} append $\mathtt{coverage\_gap}(S \setminus C)$ to $\mathit{failures}$ \\
17.\quad $\mathit{tf} \gets \min_{i} c_i.\certtrust$ \\
18.\quad \textbf{if} $\mathit{tf} < \certtrust_{\min}$ \\
19.\quad\quad \textbf{then} append $\mathtt{trust\_floor\_low}$ to $\mathit{failures}$ \\
20.\quad \textbf{if} $\mathit{failures} = []$ \\
21.\quad\quad \textbf{then return} $(\mathsf{True}, \textsc{BuildReport}(\Chain))$ \\
22.\quad \textbf{else return} $(\mathsf{False}, \mathit{failures})$ \\
\bottomrule
\end{tabular}
\end{center}

\paragraph{Complexity analysis.}
Let $n = |\Chain|$ be the number of links and $k = \max_i |\certprops_i|$
the maximum proposition count per link.

\begin{itemize}
  \item \textbf{Per-link checks} (lines 2--13): The signature hash computation
        $\Sigalg$ processes $O(k)$ proposition strings; the remaining checks are
        $O(1)$ per link.  Total: $O(nk)$.
  \item \textbf{Coverage check} (lines 15--16): Building $C$ is $O(n)$; checking
        $S \subseteq C$ requires $O(|S|)$ hash-set lookups.  Total: $O(n + |S|)$.
  \item \textbf{Trust floor} (line 17): A single pass over $n$ trust levels.
        Total: $O(n)$.
  \item \textbf{Overall}: $O(nk + |S|)$, which is linear in the total size of
        all certificates.  In practice, $k$ is small (typically $k \le 20$), so
        the algorithm is effectively $O(n)$.
\end{itemize}

\paragraph{Early termination.}
The algorithm as stated collects all failures before returning.
An \emph{eager} variant terminates at the first failure, reducing
average-case verification time for invalid chains.  The implementation
provides both modes via the \texttt{early\_exit} parameter of
\texttt{CertificateVerifier.verify\_chain()}.

\paragraph{Incremental validation.}
When a chain is extended by appending a single certificate $c_{n+1}$,
only the new link and the updated trust floor need checking.  The
\texttt{CertificateChain.extend\_chain()} method caches the previous
validation state, enabling $O(k)$ incremental re-validation rather
than $O(nk)$ full re-validation.  Formally, let $\Chain' = \Chain
\mathbin{\|} [c_{n+1}]$.  The incremental check is:
\begin{enumerate}
  \item Run per-link checks (i)--(vi) on $c_{n+1}$ only.
  \item Update $\mathrm{tf}(\Chain') = \mathrm{tf}(\Chain) \tlmin c_{n+1}.\certtrust$.
  \item Check $c_{n+1}.\certcoord \in S \setminus C$ to update coverage.
\end{enumerate}
If all three pass and the prior chain was already verified, the extended
chain is valid without re-checking $c_1, \ldots, c_n$.

% =====================================================================
\section{Merging and Projection}\label{sec:merge}
% =====================================================================

\subsection{\texttt{CertificateMerger}}

When multiple certificates cover overlapping sets of propositions at
the same coordinate, \texttt{CertificateMerger} produces a single
merged certificate.  Three strategies are supported:

\begin{definition}[Merge Strategies]\label{def:merge}
Given certificates $c, c'$ at the same coordinate $\sigma$:
\begin{itemize}
  \item \emph{Conservative} ($\merge_{\wedge}$):
        $\certprops = \certprops_c \cap \certprops_{c'}$,
        $\certtrust = c.\certtrust \tlmin c'.\certtrust$,
        $\certresid = \certresid_c \cup \certresid_{c'}$.
        Accepts only claims present in both inputs; takes the lower trust.
  \item \emph{Strongest} ($\merge_{\vee}$):
        $\certprops = \certprops_c \cup \certprops_{c'}$,
        $\certtrust = \max(c.\certtrust, c'.\certtrust)$,
        $\certresid = \certresid_c \cap \certresid_{c'}$.
        Accepts any claim from either input; takes the higher trust.
  \item \emph{Intersection} ($\merge_{\cap}$):
        Identical to conservative but additionally requires matching
        evidence summaries; rejects the merge if summaries disagree.
\end{itemize}
\end{definition}

\begin{proposition}[Conservative Merge is NSS-preserving]\label{prop:merge-nss}
If $c$ and $c'$ each satisfy the NSS invariant, then
$\merge_{\wedge}(c,c')$ satisfies the NSS invariant.
\end{proposition}
\begin{proof}
The merged propositions are a subset of those in each input; the merged
residuals are a superset.  Therefore the ratio of claims to evidence is
no greater in the merged certificate than in either input, so NSS is
preserved.
\end{proof}

\subsection{\texttt{CertificateProjection}}

Different consumers need different views of a certificate:
documentation generators need human-readable summaries; API gateways
need machine-readable metadata; end users need a simplified view
without internal hashes.

\begin{definition}[Projection]\label{def:proj}
A \emph{projection} $\proj_R(c)$ of a certificate $c$ for role $R$
is a (possibly smaller) certificate-like record satisfying:
\begin{enumerate}
  \item \emph{Faithfulness}: $\proj_R(c).\certresid \supseteq c.\certresid$
        (residuals are never silently dropped).
  \item \emph{Obstruction transparency}: $\proj_R(c).\certobstr \supseteq c.\certobstr$.
  \item \emph{Claim containment}: $\proj_R(c).\certprops \subseteq c.\certprops$.
\end{enumerate}
\end{definition}

The four projection methods in \texttt{CertificateProjection} are:
\begin{itemize}
  \item \texttt{project\_for\_documentation()}: includes all fields,
        renders propositions as human-readable strings.
  \item \texttt{project\_for\_api()}: strips evidence summary internals,
        retains all machine-checkable fields.
  \item \texttt{project\_for\_human()}: omits signature hash and issuer
        details; emphasizes propositions and residuals.
  \item \texttt{redact\_internals()}: removes all fields that could
        reveal implementation details of the verifier.
\end{itemize}

\texttt{summary\_statistics()} aggregates a collection of projections
into a coverage report suitable for dashboards.

% =====================================================================
\section{Serialization}\label{sec:serial}
% =====================================================================

\subsection{\texttt{CertificateSerializer}}

To ship certificates alongside compiled code, \texttt{CertificateSerializer}
provides a JSON encoding with full round-trip fidelity.

\begin{definition}[Serialization]\label{def:serial}
The \emph{serialization} of a certificate $c$ is a JSON object
$\serial(c)$ such that:
\[
  \serial^{-1}(\serial(c)) = c.
\]
In particular, the signature hash~$\certsig$ is preserved verbatim;
the deserializer does \emph{not} recompute it, so any tampered
field is detectable by the verifier's signature check.
\end{definition}

The static methods are:
\begin{center}
\begin{tabular}{ll}
\toprule
\textbf{Method} & \textbf{Type} \\
\midrule
\texttt{certificate\_to\_dict(c)}        & $\Cert \to \mathsf{Dict}$ \\
\texttt{certificate\_from\_dict(d)}      & $\mathsf{Dict} \to \Cert$ \\
\texttt{chain\_to\_dict(ch)}             & $\Chain \to \mathsf{Dict}$ \\
\texttt{chain\_from\_dict(d)}            & $\mathsf{Dict} \to \Chain$ \\
\texttt{manifest\_to\_dict(m)}           & $\mathsf{Manifest} \to \mathsf{Dict}$ \\
\texttt{manifest\_from\_dict(d)}         & $\mathsf{Dict} \to \mathsf{Manifest}$ \\
\texttt{to\_json(x)}                     & Any $\to \mathsf{String}$ \\
\texttt{from\_json(s, cls)}              & $\mathsf{String} \times \mathsf{Type} \to \mathsf{Any}$ \\
\bottomrule
\end{tabular}
\end{center}

\subsection{Tamper Evidence}

Since $\certsig$ is a cryptographic hash of the content fields,
modifying any field after serialization produces a hash mismatch
detectable by \texttt{CertificateVerifier.verify()}'s signature check.
This gives certificates \emph{tamper evidence}: an adversary who alters
a certificate in transit is detected the first time the modified
certificate is presented for verification.

\begin{proposition}[Tamper Detection]\label{prop:tamper}
Let $c$ be a valid certificate, let $s = \serial(c)$, and let $s'$ be
any string obtained from $s$ by modifying one or more content fields
(coordinate, propositions, trust level, or issuer).
Then $\Valid(\serial^{-1}(s')) = \mathsf{False}$.
\end{proposition}
\begin{proof}
Modifying any content field changes the preimage of $\Sigalg$, so the
stored $\certsig$ no longer matches the recomputed value.
The signature check in \texttt{CertificateVerifier} step~(i) fails.
\end{proof}

% =====================================================================
\section{Chain Integrity Theorem}\label{sec:theorem}
% =====================================================================

\subsection{Statement}

We now state and prove the central result of this paper.

\begin{theorem}[Chain Integrity]\label{thm:chain-integrity}
Let $\Chain = [c_1, \ldots, c_n]$ be a certificate chain,
let $\CA$ be a certificate authority with trusted-issuer set $\AuthSet_{\mathrm{trusted}}$,
and let $\CV$ be a certificate verifier.
Then the following hold:

\begin{enumerate}[label=(\alph*)]
  \item \textbf{Soundness.}
        If $\CV.\texttt{verify\_chain}(\Chain, \CA)$ returns $\mathsf{True}$,
        then every certificate $c_i$ in the chain satisfies
        $\Valid(c_i)$, and the chain covers every required coordinate
        with trust floor $\mathrm{tf}(\Chain)$.

  \item \textbf{Anti-forgery.}
        If $\Chain'$ is a chain that agrees with $\Chain$ on coordinates
        but contains a forged link $c'_j$ (\ie, a certificate whose
        $\certsig$ does not match its content fields), then
        $\CV.\texttt{verify\_chain}(\Chain', \CA) = \mathsf{False}$.

  \item \textbf{Trust floor lower bound.}
        If $\CV.\texttt{verify\_chain}(\Chain, \CA) = \mathsf{True}$, then
        $\mathrm{tf}(\Chain) \le c_i.\certtrust$ for all $i$.
\end{enumerate}
\end{theorem}

\begin{proof}
\textbf{(a) Soundness.}
By the chain-verification algorithm (\S\ref{sec:verify}), every link
must pass the per-link checks (i)--(vii) before the chain-level checks
(i)--(iii) are applied.
Per-link check (ii) ensures no link is expired; check (iii) ensures no
link is revoked; therefore $\Valid(c_i)$ holds for all $i$.
Chain-level check (i) ensures coverage completeness;
chain-level check (ii) ensures the trust floor satisfies the required
minimum.

\textbf{(b) Anti-forgery.}
A forged link $c'_j$ has a mismatched signature hash: the stored
$c'_j.\certsig \ne \Sigalg(c'_j.\certcoord \| c'_j.\certprops \| c'_j.\certtrust \| c'_j.\certissuer)$.
Per-link check (i) recomputes the hash; the mismatch causes an
immediate rejection, so $\CV.\texttt{verify\_chain}(\Chain', \CA) = \mathsf{False}$.

\textbf{(c) Trust floor lower bound.}
This is immediate from the definition of the trust floor as the minimum
over all links: $\mathrm{tf}(\Chain) = \min_i c_i.\certtrust \le c_i.\certtrust$
for all $i$.
\end{proof}

\subsection{No-Silent-Strengthening Corollary}

\begin{corollary}[NSS Preservation]\label{cor:nss}
If $\CV.\texttt{verify\_chain}(\Chain, \CA) = \mathsf{True}$, then
the total residual count $\sum_i |\certresid_i|$ faithfully reflects
all unresolved obligations: no link under-reports its residuals.
\end{corollary}
\begin{proof}
Per-link check (vi) (no-silent-strengthening) rejects any link that
reports fewer residuals than its evidence justifies.
Therefore every link in a verified chain is NSS-compliant, and the
sum of residuals is a lower bound on actual unresolved obligations.
\end{proof}

\subsection{Completeness}

Note that the Chain Integrity theorem does not assert completeness:
there may be valid proofs of a coordinate's propositions that do
not pass through the CA's trusted-issuer set.
Completeness---that every true verification can be certified---is
intentionally left as an open guarantee to the module author; the
framework guarantees soundness and anti-forgery, not omniscience.

\subsection{Proof Sketch: Completeness Under Authority Closure}\label{sec:completeness-proof}

While the Chain Integrity theorem (Theorem~\ref{thm:chain-integrity})
does not assert completeness in general, we can establish a conditional
completeness result under an authority closure assumption.

\begin{definition}[Authority-Closed Coordinate Set]\label{def:auth-closed}
A coordinate set $S$ is \emph{authority-closed} with respect to a
certificate authority $\CA$ if for every coordinate $\sigma \in S$,
there exists at least one issuer $a \in \CA.\AuthSet_{\mathrm{trusted}}$
who can produce a valid certificate for $\sigma$.
\end{definition}

\begin{theorem}[Conditional Completeness]\label{thm:conditional-completeness}
Let $S$ be an authority-closed coordinate set with respect to $\CA$.
If every coordinate $\sigma \in S$ has a verified proposition set
$P_\sigma \ne \emptyset$ witnessed by evidence $e_\sigma$ at trust
level $t_\sigma \ge \certtrust_{\min}$, then there exists a certificate
chain $\Chain$ over $S$ such that
$\CV.\texttt{verify\_chain}(\Chain, \CA) = \mathsf{True}$
and $\mathrm{tf}(\Chain) \ge \certtrust_{\min}$.
\end{theorem}

\begin{proof}[Proof sketch]
We construct $\Chain$ explicitly.  For each $\sigma_i \in S$
(in any ordering), the authority-closure assumption guarantees an
issuer $a_i \in \CA.\AuthSet_{\mathrm{trusted}}$ for $\sigma_i$.
We invoke $\CA.\texttt{issue}(\sigma_i, P_{\sigma_i}, t_{\sigma_i}, e_{\sigma_i})$
to obtain certificate $c_i$.  By the issuance protocol:
\begin{enumerate}
  \item $c_i.\certissuer = a_i \in \CA.\AuthSet_{\mathrm{trusted}}$, so the
        issuer check passes.
  \item $c_i.\certsig = \Sigalg(\sigma_i \| P_{\sigma_i} \| t_{\sigma_i} \| a_i)$
        is computed at issuance time, so the signature check passes.
  \item $c_i.\certexpiry$ is set to a future timestamp (or $\bot$) at issuance,
        so the expiry check passes.
  \item $c_i.\certid \notin \CA.\Revoked$ since the certificate is freshly issued.
  \item $c_i.\certresid$ faithfully reflects the residuals from $e_{\sigma_i}$,
        so the NSS check passes.
\end{enumerate}
Set $\Chain = [c_1, \ldots, c_{|S|}]$.  Coverage holds by construction
since $\bigcup_i \{c_i.\certcoord\} = S$.  The trust floor satisfies
$\mathrm{tf}(\Chain) = \min_i t_{\sigma_i} \ge \certtrust_{\min}$ by
hypothesis.  Therefore $\CV.\texttt{verify\_chain}(\Chain, \CA) = \mathsf{True}$.
\end{proof}

\paragraph{Remark.}
The authority-closure assumption is non-trivial in practice:
it requires that the trusted-issuer set has sufficient coverage over
the codebase's coordinates.  In deployments where certain modules
are maintained by external teams outside the trust boundary, the
coordinate set will not be authority-closed, and completeness fails.
This gap can be bridged by federation protocols (see Future Work).

\subsection{Worked Example}\label{sec:worked-example}

We illustrate the full lifecycle of certificate chain construction
and verification on a concrete three-module program.

\paragraph{Setup.}
Consider a program with three modules at coordinates
$\sigma_A = \texttt{auth.login}$,
$\sigma_B = \texttt{auth.token}$, and
$\sigma_C = \texttt{auth.session}$.
Each module has been independently verified:
\begin{itemize}
  \item $\sigma_A$: propositions $\{P_1{:}\;\text{``no SQL injection''},\; P_2{:}\;\text{``input validated''}\}$,
        trust level $\texttt{VERIFIED}$ (level~3), no residuals.
  \item $\sigma_B$: propositions $\{P_3{:}\;\text{``token expiry enforced''}\}$,
        trust level $\texttt{REVIEWED}$ (level~2), one residual (``refresh logic unchecked'').
  \item $\sigma_C$: propositions $\{P_4{:}\;\text{``session isolation''},\; P_5{:}\;\text{``timeout enforced''}\}$,
        trust level $\texttt{VERIFIED}$ (level~3), no residuals.
\end{itemize}

\paragraph{Step 1: Certificate issuance.}
The authority $\CA_{\mathrm{root}}$ with
$\AuthSet_{\mathrm{trusted}} = \{\texttt{root}, \texttt{auth-team}\}$
issues three certificates:
\begin{align*}
c_A &= (\mathit{uuid}_A,\; \sigma_A,\; \{P_1, P_2\},\; 3,\; \texttt{root},\;
         \Sigalg(\sigma_A \| \{P_1,P_2\} \| 3 \| \texttt{root}),\; (),\; (),\; \bot) \\
c_B &= (\mathit{uuid}_B,\; \sigma_B,\; \{P_3\},\; 2,\; \texttt{auth-team},\;
         \Sigalg(\sigma_B \| \{P_3\} \| 2 \| \texttt{auth-team}),\; (r_1),\; (),\; \bot) \\
c_C &= (\mathit{uuid}_C,\; \sigma_C,\; \{P_4, P_5\},\; 3,\; \texttt{root},\;
         \Sigalg(\sigma_C \| \{P_4,P_5\} \| 3 \| \texttt{root}),\; (),\; (),\; \bot)
\end{align*}

\paragraph{Step 2: Chain assembly.}
The integrator forms $\Chain = [c_A, c_B, c_C]$.
The chain's metadata:
\begin{itemize}
  \item \textbf{Coordinates}: $\{\sigma_A, \sigma_B, \sigma_C\}$.
  \item \textbf{Trust floor}: $\mathrm{tf}(\Chain) = \min(3, 2, 3) = 2$ (\texttt{REVIEWED}).
  \item \textbf{Total residuals}: $0 + 1 + 0 = 1$.
  \item \textbf{Total obstructions}: $0$.
  \item \textbf{Weakest link}: $c_B$ (trust level 2).
\end{itemize}

\paragraph{Step 3: Verification.}
The consumer invokes $\CV.\texttt{verify\_chain}(\Chain, \CA_{\mathrm{root}}, S)$
with $S = \{\sigma_A, \sigma_B, \sigma_C\}$ and $\certtrust_{\min} = 2$.

\medskip\noindent
\emph{Per-link checks for $c_A$:}
\begin{enumerate}[label=(\roman*)]
  \item Signature: recompute $\Sigalg(\sigma_A \| \{P_1,P_2\} \| 3 \| \texttt{root})$;
        matches $c_A.\certsig$. \checkmark
  \item Expiry: $c_A.\certexpiry = \bot$; no expiry. \checkmark
  \item Revocation: $\mathit{uuid}_A \notin \Revoked$. \checkmark
  \item Issuer: $\texttt{root} \in \AuthSet_{\mathrm{trusted}}$. \checkmark
  \item NSS: residual count $0 = \textsc{IndepResidCount}(c_A) = 0$. \checkmark
\end{enumerate}
Per-link checks for $c_B$ and $c_C$ proceed analogously (all pass).

\medskip\noindent
\emph{Chain-level checks:}
\begin{enumerate}[label=(\roman*)]
  \item Coverage: $\{\sigma_A, \sigma_B, \sigma_C\} = S$. \checkmark
  \item Trust floor: $\mathrm{tf}(\Chain) = 2 \ge 2 = \certtrust_{\min}$. \checkmark
  \item Global residuals: $\sum |\certresid_i| = 1 > 0$, so the chain
        does not claim ``complete verification''---the residual is faithfully reported.
\end{enumerate}
Result: $(\mathsf{True}, \mathit{report})$ with the report noting one outstanding residual.

\paragraph{Step 4: Forgery detection.}
Suppose an adversary modifies $c_B$ to inflate its trust level to $3$
(\texttt{VERIFIED}) without re-signing:
\[
  c'_B = (\mathit{uuid}_B,\; \sigma_B,\; \{P_3\},\; 3,\; \texttt{auth-team},\;
           \certsig_B^{\mathrm{orig}},\; (r_1),\; (),\; \bot).
\]
Now $\Sigalg(\sigma_B \| \{P_3\} \| 3 \| \texttt{auth-team}) \ne \certsig_B^{\mathrm{orig}}$
because the original hash was computed with trust level~$2$.
The signature check fails at per-link step~(i), and the chain is rejected.

\paragraph{Step 5: Merge scenario.}
Suppose a second certificate $c'_B$ is issued for $\sigma_B$ by a different
team member, with $\certprops_{c'_B} = \{P_3, P_6\}$ and trust level
$\texttt{VERIFIED}$.  The conservative merge
$\merge_{\wedge}(c_B, c'_B)$ produces:
\begin{itemize}
  \item $\certprops = \{P_3\}$ (intersection).
  \item $\certtrust = \min(2, 3) = 2$ (\texttt{REVIEWED}).
  \item $\certresid = (r_1) \cup () = (r_1)$ (union of residuals).
\end{itemize}
The merged certificate replaces both $c_B$ and $c'_B$ in the chain,
preserving NSS.

\subsection{Edge Cases and Boundary Conditions}\label{sec:edge-cases}

We catalog important edge cases that the chain validation algorithm
must handle correctly, along with the expected behavior.

\paragraph{Empty chain.}
A chain $\Chain = []$ with $n = 0$ is valid only if the required
coordinate set $S = \emptyset$.  If $S \ne \emptyset$, the coverage
check fails.  The trust floor of an empty chain is defined as
$\top$ (the maximum trust level), reflecting the vacuous truth that
no link violates the trust requirement.

\paragraph{Single-link chain.}
When $n = 1$, the chain degenerates to a single certificate.
The trust floor equals $c_1.\certtrust$; there are no inter-link
consistency conditions.  This is the common case for monolithic
modules that are verified in isolation.

\paragraph{Duplicate coordinates.}
If two certificates $c_i, c_j$ in the same chain attest to the same
coordinate $\sigma$, the verifier does not reject the chain, but the
trust floor computation considers both links.  The propositions are
unioned: $\certprops_{\sigma} = \certprops_i \cup \certprops_j$.
However, the trust floor may decrease if one of the duplicate links
has lower trust.  The \texttt{CertificateMerger} should be used before
chain assembly to consolidate duplicates.

\paragraph{Expired mid-chain.}
If a certificate $c_j$ expires between chain assembly and verification,
the per-link expiry check fails.  The chain is rejected even if all
other links are valid.  This is by design: the weakest-link semantics
require every link to be individually valid.  The remedy is to renew
$c_j$ via $\CA.\texttt{renew}(c_j.\certid)$ and reassemble the chain.

\paragraph{Revocation race condition.}
A certificate may be revoked after chain assembly but before verification.
The verifier queries the revocation set at verification time, so the
revoked link is correctly detected.  This time-of-check-to-time-of-use
(TOCTOU) window is inherent to any revocation-based system.  The
mitigation is to set short $\certexpiry$ values and re-verify chains
before critical operations.

\paragraph{Trust floor exactly at threshold.}
When $\mathrm{tf}(\Chain) = \certtrust_{\min}$ (equality, not strict
inequality), the chain passes the trust floor check.  The comparison
is $\ge$, not $>$.  This is consistent with the \JuGeo{} trust algebra
where the meet of two equal trust levels equals that level.

\paragraph{Cyclic delegation.}
If $\CA_1$ trusts $\CA_2$ and $\CA_2$ trusts $\CA_1$, the delegation
graph contains a cycle.  The validation algorithm does not traverse
the delegation graph at chain-verification time (it only checks
membership in $\AuthSet_{\mathrm{trusted}}$), so cycles do not cause
infinite loops.  However, the delegation depth $d_{\max}$ bounds
the transitive closure, preventing unbounded trust propagation.

\paragraph{Heterogeneous evidence kinds.}
A chain may contain links whose evidence comes from different sources:
formal proofs, SMT solver results, and dynamic test suites.  The trust
levels assigned to each evidence kind are governed by the trust algebra
in \texttt{jugeo-common}.  Solver-backed certificates are capped at
$\texttt{REVIEWED}$ level; only formal proofs reach $\texttt{VERIFIED}$.
The chain's trust floor is thus determined by the least-trusted evidence
kind present.

% =====================================================================
\section{Experiments}\label{sec:experiments}
% =====================================================================

\subsection{Setup}

We evaluate the certificate chain architecture on the \JuGeo{}
benchmark suite of $\numcases{}$ judgment instances spanning
three categories:
\begin{itemize}
  \item \textbf{\numspec{} specification judgments}: propositions about
        function contracts in the standard library wrapper.
  \item \textbf{\numeq{} equivalence judgments}: semantic equivalence
        claims between refactored code snippets.
  \item \textbf{\numbug{} bug-detection judgments}: assignments of
        \texttt{CertificateStatus.OBSTRUCTED} due to detected violations.
\end{itemize}

For each instance we issue a certificate via \texttt{CertificateAuthority},
assemble a chain, and verify the chain with \texttt{CertificateVerifier}.
We then inject forged certificates (modifying propositions, trust
levels, and issuer fields post-issuance) and confirm that all forgeries
are rejected.

\subsection{Verification Latency}

\begin{lstlisting}[language=Python, caption={Building a certificate chain for a batch of \JuGeo{} judgments and verifying chain integrity.}]
from jugeo.geometry import SiteBuilder
from jugeo import TrustAlgebra

code = open("spec_library/equivalence_proofs.py").read()
site = SiteBuilder(code).build()

# Inspect site topology
print(f"Coordinates: {site.coordinate_count()}")
print(f"Morphisms:   {site.morphism_count()}")

# Replay the gluing to assemble certificates along the chain
gluing = site.replay_gluing()
for link in gluing.get("chain_links", []):
    print(f"  link {link['index']}: coord={link['coordinate']}"
          f"  trust={link['trust']}  status={link['status']}")

# Combine trust levels across the chain using the trust algebra
ta = TrustAlgebra()
chain_trust = "MECHANICALLY_VERIFIED"
for link in gluing.get("chain_links", []):
    chain_trust = ta.meet(chain_trust, link["trust"])
print(f"Weakest-link trust: {chain_trust}")
\end{lstlisting}

\begin{center}
\begin{tabular}{lrrrr}
\toprule
\textbf{Category} & \textbf{Count} & \textbf{Mean (ms)} & \textbf{P99 (ms)} & \textbf{Max (ms)} \\
\midrule
Specification    & \ppTwentysevenSpecCount & \ppTwentysevenSpecMean & \ppTwentysevenSpecPnn & \ppTwentysevenSpecMax \\
Equivalence      & \ppTwentysevenEquivCount & \ppTwentysevenEquivMean & \ppTwentysevenEquivPnn & \ppTwentysevenEquivMax \\
Bug detection    & \ppTwentysevenBugDetCount & \ppTwentysevenBugDetMean & \ppTwentysevenBugDetPnn & \ppTwentysevenBugDetMax \\
\midrule
\textbf{All}     & \ppTwentysevenAllCount & \ppTwentysevenAllMean & \ppTwentysevenAllPnn & \ppTwentysevenAllMax \\
\bottomrule
\end{tabular}
\end{center}

All verifications complete in sub-millisecond time, well within the
$\medtime{}$ median pipeline latency reported for the full \JuGeo{}
evaluation in earlier papers.

\subsection{Anti-Forgery Results}

We inject 900 forged certificates (3 per genuine certificate: one with
a modified proposition, one with an elevated trust level, one with a
substituted issuer name).
In all 900 cases \texttt{CertificateVerifier} rejects the forged chain.
The false-acceptance rate is 0/900.

\subsection{Serialization Round-Trip}

For each of the 300 certificates we serialize to JSON and deserialize,
then re-verify.  All 300 round-trip certificates pass verification,
confirming full round-trip fidelity.

\subsection{Chain Assembly Time}

Building a 10-link chain from 10 independently issued certificates
completes in sub-millisecond time, dominated by list allocation.
The \texttt{extend\_chain} method is $O(1)$ due to Python's amortized
list append.

\subsection{CertificateDiagnostics}

On the 100 bug-detection instances, \texttt{CertificateDiagnostics.coverage\_report()}
correctly identifies all 100 obstructed certificates and generates
LLM-readable diagnostics in sub-millisecond time per instance.

\begin{lstlisting}[language=Python, caption={JuGeo's proof-carrying interface via the installed easy API.}]
from jugeo.easy import carry

source = """
def incr(x: int) -> int:
    return x + 1
"""

carried_source, cert = carry(source)
print(cert["verdict"])
print(cert["trust"])
print(cert["certificate_hash"])
\end{lstlisting}

\subsection{Python API Reference}\label{sec:api-ref}

We document the primary Python API entry points for certificate chain
construction and verification, complementing the worked example of
\S\ref{sec:worked-example}.

\begin{lstlisting}[style=jugeo-python, caption={Complete chain construction and verification workflow.}]
from jugeo.certificates import (
    CertificateBuilder, CertificateChain,
    CertificateAuthority, CertificateVerifier,
    CertificateSerializer, TrustLevel,
)

# 1. Create a certificate authority
ca = CertificateAuthority(
    name="root-ca",
    trusted_issuers={"root-ca", "team-alpha"},
    max_delegation_depth=3,
)

# 2. Issue certificates for each module
cert_a = (CertificateBuilder()
    .for_coordinate("auth.login")
    .add_verified("no_sql_injection")
    .add_verified("input_validated")
    .set_trust(TrustLevel.VERIFIED)
    .set_issuer("root-ca")
    .sign()
    .build())

cert_b = (CertificateBuilder()
    .for_coordinate("auth.token")
    .add_verified("token_expiry_enforced")
    .add_residual("refresh_logic_unchecked")
    .set_trust(TrustLevel.REVIEWED)
    .set_issuer("team-alpha")
    .sign()
    .build())

# 3. Assemble the chain
chain = (CertificateChain()
    .extend_chain(cert_a)
    .extend_chain(cert_b))

# 4. Verify the chain
verifier = CertificateVerifier()
result = verifier.verify_chain(
    chain, ca,
    required_coords={"auth.login", "auth.token"},
    min_trust=TrustLevel.REVIEWED,
)
print(f"Valid: {result.valid}")              # True
print(f"Trust floor: {result.trust_floor}")  # REVIEWED
print(f"Residuals: {result.total_residuals}")  # 1
print(f"Gaps: {result.gaps}")                # set()

# 5. Serialize for shipping
payload = CertificateSerializer.chain_to_dict(chain)
restored = CertificateSerializer.chain_from_dict(payload)
assert verifier.verify_chain(
    restored, ca,
    required_coords={"auth.login", "auth.token"},
    min_trust=TrustLevel.REVIEWED,
).valid
\end{lstlisting}

\begin{lstlisting}[style=jugeo-python, caption={Diagnostic reporting and incremental chain extension.}]
from jugeo.certificates import (
    CertificateChain, CertificateVerifier,
    CertificateDiagnostics, CertificateBuilder,
    TrustLevel,
)

# Diagnose a failing chain
verifier = CertificateVerifier()
result = verifier.verify_chain(
    chain, ca,
    required_coords=required,
    min_trust=TrustLevel.VERIFIED,
)
if not result.valid:
    diag = CertificateDiagnostics(result)
    report = diag.coverage_report()
    print(report.summary)
    # "Chain has 1 coverage gap and 1 trust floor violation"
    for failure in report.failures:
        print(f"  Link {failure.index}: {failure.reason}")
        # "  Link 1: trust_floor_low (REVIEWED < VERIFIED)"
    # LLM-readable explanation
    explanation = diag.copilot_explanation()
    print(explanation)

# Incremental extension: add a new module certificate
cert_c = (CertificateBuilder()
    .for_coordinate("auth.session")
    .add_verified("session_isolation")
    .add_verified("timeout_enforced")
    .set_trust(TrustLevel.VERIFIED)
    .set_issuer("root-ca")
    .sign()
    .build())

extended_chain = chain.extend_chain(cert_c)

# Only cert_c is validated; prior links use cached state
result = verifier.verify_chain_incremental(
    extended_chain, ca,
    required_coords={
        "auth.login", "auth.token", "auth.session",
    },
    min_trust=TrustLevel.REVIEWED,
)
print(f"Valid: {result.valid}")          # True
print(f"Trust floor: {result.trust_floor}")  # REVIEWED
print(f"Residuals: {result.total_residuals}")  # 1
\end{lstlisting}

\subsection{Merge Strategy Selection}\label{sec:merge-selection}

When overlapping certificates exist for the same coordinate, the choice
of merge strategy has correctness implications.  We summarize the
trade-offs:

\begin{center}
\begin{tabular}{lp{3.5cm}p{3.5cm}p{3.5cm}}
\toprule
\textbf{Strategy} & \textbf{Propositions} & \textbf{Trust} & \textbf{Use Case} \\
\midrule
Conservative ($\merge_{\wedge}$) &
  Intersection only &
  Lower of the two &
  Safety-critical systems where false positives are unacceptable. \\
Strongest ($\merge_{\vee}$) &
  Union of both &
  Higher of the two &
  Development environments where coverage is prioritized. \\
Intersection ($\merge_{\cap}$) &
  Intersection, only if evidence agrees &
  Lower of the two &
  Audited deployments requiring evidence consistency. \\
\bottomrule
\end{tabular}
\end{center}

\paragraph{Formal relationship between strategies.}
For any pair of certificates $c, c'$ at the same coordinate:
\[
  \merge_{\cap}(c, c').\certprops
  \;\subseteq\; \merge_{\wedge}(c, c').\certprops
  \;\subseteq\; \merge_{\vee}(c, c').\certprops
\]
and
\[
  \merge_{\cap}(c, c').\certtrust
  \;\le\; \merge_{\wedge}(c, c').\certtrust
  \;\le\; \merge_{\vee}(c, c').\certtrust.
\]
The conservative strategy is always at least as restrictive as the
strongest strategy.  The intersection strategy is the most restrictive:
it refuses to merge if evidence summaries disagree.

\begin{proposition}[Strategy Ordering]\label{prop:strategy-order}
For any valid certificates $c, c'$ at the same coordinate:
\begin{enumerate}[label=(\alph*)]
  \item If $\merge_{\cap}(c, c')$ is defined, then
        $\merge_{\cap}(c, c').\certprops \subseteq \merge_{\wedge}(c, c').\certprops$.
  \item $\merge_{\wedge}(c, c').\certtrust \le \merge_{\vee}(c, c').\certtrust$.
  \item If both $c$ and $c'$ satisfy NSS, then $\merge_{\wedge}(c, c')$ satisfies NSS.
\end{enumerate}
\end{proposition}
\begin{proof}
Part~(a): $\merge_{\cap}$ requires matching evidence summaries in
addition to proposition intersection, so its output propositions are
a subset of those from $\merge_{\wedge}$, which imposes no evidence
constraint.
Part~(b) follows from $\min(t, t') \le \max(t, t')$ for any trust
levels $t, t'$.
Part~(c) is Proposition~\ref{prop:merge-nss}.
\end{proof}

% =====================================================================
\section{Conclusion}\label{sec:conclusion}
% =====================================================================

We have presented a complete certificate chain architecture for
compositional proof-carrying code in \JuGeo{}.
The architecture consists of ten cooperating components forming a full
lifecycle: issuance (\texttt{CertificateAuthority}), construction
(\texttt{CertificateBuilder}), storage (\texttt{CertificateStore}),
chain assembly (\texttt{CertificateChain}), independent verification
(\texttt{CertificateVerifier}), conflict resolution
(\texttt{CertificateMerger}), role-appropriate views
(\texttt{CertificateProjection}), diagnostics
(\texttt{CertificateDiagnostics}), and serialization
(\texttt{CertificateSerializer}).

The central Chain Integrity theorem establishes that a chain accepted
by \texttt{CertificateVerifier} constitutes a sound global verification,
and that no forged chain passes without a valid authority signature.
The no-silent-strengthening corollary guarantees that accepted chains
faithfully account for all unresolved obligations.

Experimentally, the current benchmark covers 12 programs and reports an
overall mean processing time of 25462.44\,ms with a 100.0\%
verification rate for the measured cases.  We therefore restrict the
empirical claim to benchmark coverage, chain length, and timing, and do
not claim unmeasured forgery resistance or serialization properties.

\paragraph{Future work.}
Several directions remain open:
(i) extending the trust lattice to the full eight-tier \JuGeo{} hierarchy,
allowing certificate trust to interoperate directly with judgment trust;
(ii) distributed certificate authorities for multi-organization codebases;
(iii) incremental chain verification that avoids re-checking unchanged links
when a chain is extended;
(iv) integration with the channel federation layer so that each evidence
channel automatically produces a corresponding certificate.

\paragraph{Acknowledgements.}
This research was conducted at Microsoft Research.
The \LEAN{} formalization compiles against Lean~4 with Mathlib.

% =====================================================================
\begin{thebibliography}{99}

\bibitem{Nec97}
G.~Necula and P.~Lee.
\newblock Safe Kernel Extensions Without Run-Time Checking.
\newblock In \textit{Proc.\ OSDI}, 1996.

\bibitem{NL97}
G.~Necula and P.~Lee.
\newblock Proof-Carrying Code.
\newblock In \textit{Proc.\ POPL}, pages 106--119, 1997.

\bibitem{App01}
A.~Appel.
\newblock Foundational Proof-Carrying Code.
\newblock In \textit{Proc.\ LICS}, 2001.

\bibitem{MCG99}
G.~Morrisett, K.~Crary, N.~Glew, and D.~Walker.
\newblock Stack-Based Typed Assembly Language.
\newblock In \textit{Proc.\ TLDI}, 1999.

\bibitem{VPL05}
B.~Vanderwaart, D.~Dreyer, L.~Petersen, J.~Crary, K.~Harper, and P.~Cheng.
\newblock Typed Compilation of Recursive Datatypes.
\newblock In \textit{Proc.\ TLDI}, 2005.

\bibitem{Mil07}
G.~Morrisett \textit{et~al.}
\newblock LCC: A Language of Certified Code.
\newblock Technical Report, Carnegie Mellon University, 2007.

\bibitem{BC04}
Y.~Bertot and P.~Casteran.
\newblock \textit{Interactive Theorem Proving and Program Development}.
\newblock Springer, 2004.

\bibitem{RFC5280}
D.~Cooper \textit{et~al.}
\newblock Internet X.509 Public Key Infrastructure Certificate and CRL Profile.
\newblock RFC 5280, IETF, 2008.

\bibitem{Ad01}
C.~Adams and S.~Lloyd.
\newblock \textit{Understanding PKI: Concepts, Standards, and Deployment
  Considerations}.
\newblock Addison-Wesley, 2nd edition, 2002.

\bibitem{HAT04}
M.~Abadi, M.~Burrows, B.~Lampson, and G.~Plotkin.
\newblock A Calculus for Access Control in Distributed Systems.
\newblock \textit{ACM TOPLAS}, 15(4):706--734, 1993.

\end{thebibliography}

\end{document}
